\clearpage\thispagestyle{empty}
\pdfbookmark[0]{Quelle zum Brief an von Lindenau vom 14. August 1809}{lindenau1809aug14-sources}
\begin{center}\textbf{Quelle zum Brief an von Lindenau vom 14. August 1809}\end{center}
\noindent Diese KI-Übertragung ins heutige Deutsch folgt dem Brief von Carl Friedrich Gauß an Bernhard August von Lindenau in \href{https://gdz.sub.uni-goettingen.de/id/PPN236007467}{Gauß’ \textit{Werken}, Bd.~VI} (Göttingen, 1874), S.~314--315. Der Erstabdruck erschien als Beitrag X, \textit{Auszug aus einem Schreiben des Hrn. Prof. Gauß}, im Juliheft 1809 der \textit{Monatlichen Correspondenz}, Bd.~XX, S.~\href{https://api.digitale-sammlungen.de/iiif/image/v2/bsb10538610_00084/full/1600,/0/default.jpg}{78}, \href{https://api.digitale-sammlungen.de/iiif/image/v2/bsb10538610_00085/full/1600,/0/default.jpg}{79}. Die Datierungen von Heft und Brief bleiben erhalten. Das digitalisierte Exemplar stammt aus der Bayerischen Staatsbibliothek.
\par\smallskip\noindent Zum Äquatorialsektor in Mailand siehe den \href{https://www.lombardiabeniculturali.it/scienza-tecnologia/schede/ST060-00001/}{Katalogeintrag ST060-00001 von Lombardia Beni Culturali} und den \href{https://artsandculture.google.com/asset/equatorial-sector/OwGZAfaAaOv35g}{englischen Objekteintrag des Museums}. KI-Übertragung und KI-Satz.
\clearpage
